\chapter{ФИЗИЧЕСКИ МОТИВИРОВАННАЯ СУРРОГАТНАЯ МОДЕЛЬ}

\section{Выбор признаков на основе scaling \texorpdfstring{$1/a^2$}{1/a^2}}

В предыдущей главе было показано, что для рассматриваемых суперэллиптических квантовых точек величина $E_{\mathrm{kin}}=E_0+4$ приблизительно масштабируется как $1/a^2$ в выбранном design window. Это согласуется с низкоэнергетическим квадратичным пределом square-lattice tight-binding модели, где характерное волновое число ограниченного состояния масштабируется как $k \sim 1/a$, а энергия, отсчитанная от дна зоны, как $k^2$.

Поэтому признаки суррогатной модели выбираются не путём слепого перебора, а из физической структуры задачи. Для фиксированного класса формы $n$ используются геометрические параметры $a$ и $r_{\mathrm{AR}}$, где $b=a r_{\mathrm{AR}}$. В анизотропном случае второе геометрическое направление также влияет на спектр, поэтому один только признак $1/a^2$ недостаточен для описания всей сетки параметров.

В качестве физически мотивированного набора признаков используется
\begin{equation}
  x_1 = \frac{1}{a^2}, \qquad
  x_2 = \frac{1}{a^2 r_{\mathrm{AR}}}, \qquad
  x_3 = r_{\mathrm{AR}} .
\end{equation}
Первый признак кодирует основной размерный масштаб, ожидаемый из низкоэнергетического закона $E_{\mathrm{kin}}\sim 1/a^2$. Второй признак является простой формой учёта второго геометрического масштаба, связанного с $b=a r_{\mathrm{AR}}$, и по смыслу близок к inverse-area-like зависимости. Третий признак оставляет в модели прямую информацию об анизотропии формы, которая может не полностью сводиться к степенным масштабам.

Такой выбор признаков не означает, что данный набор является единственным возможным или глобально оптимальным. Он означает только, что в проверенном диапазоне параметров признаки согласованы с физической картиной низкоэнергетического confinement и могут служить компактным описанием основной зависимости спектра от размера и формы. Обобщение за пределы диапазонов $a \in [24,36]$ и $r_{\mathrm{AR}} \in [0.67,1.0]$ в данной работе не утверждается.

Целевыми величинами для основной регрессии являются $E_0$ и $dE_1=E_1-E_0$. Величина $dE_2$ сохраняется как диагностическая, но не используется как основная цель суррогатной модели из-за чувствительности к вырождению и порядку уровней, описанной в главе 3.

\section{Ridge-регрессия как интерпретируемый baseline}

Основной baseline в работе построен на Ridge-регрессии. Это линейная модель в выбранном пространстве признаков:
\begin{equation}
  \hat{y}
  =
  \beta_0
  +
  \sum_{j=1}^{3} \beta_j x_j ,
\end{equation}
где $\hat{y}$ обозначает предсказание для $E_0$ или $dE_1$, а $x_j$ --- физически мотивированные признаки. Ridge-регрессия минимизирует ошибку с добавлением $L_2$-регуляризации коэффициентов:
\begin{equation}
  \mathcal{L}
  =
  \sum_i \left(y_i-\hat{y}_i\right)^2
  +
  \alpha \sum_j \beta_j^2 .
\end{equation}

Такой baseline не является формальной игрушечной моделью. Для малых структурированных датасетов низкая ёмкость модели является преимуществом: она снижает риск подгонки к отдельным точкам сетки и делает результат устойчивее к небольшим lattice-эффектам. $L_2$-регуляризация стабилизирует коэффициенты, а линейность в физически выбранных признаках облегчает интерпретацию. Если такая модель хорошо аппроксимирует спектральную характеристику, это означает, что существенная часть зависимости уже описывается проверенным scaling и простой анизотропной поправкой.

Перед обучением признаки стандартизируются, чтобы регуляризация действовала сопоставимо на разные компоненты признакового вектора. При этом метрики качества вычисляются в исходных физических единицах hopping-параметра $t$, а не в нормированных единицах. Параметры baseline фиксируются заранее; tuning гиперпараметров не используется как дополнительный способ улучшить результат после просмотра тестовых ошибок.

В данной методологии Ridge остаётся предпочтительной моделью по умолчанию, пока более сложная модель не демонстрирует устойчивого и заранее определённого преимущества. Такой подход соответствует физической задаче: цель состоит не в максимизации сложности surrogate-модели, а в получении проверяемой аппроксимации спектра модельной наноструктуры.

\section{Протоколы структурированной валидации LOAO и LOARO}

Обычное случайное разбиение train/test не является подходящим основным протоколом для данного датасета. Геометрии образуют регулярную параметрическую сетку по $a$ и $r_{\mathrm{AR}}$, поэтому при random split в обучающей и тестовой выборках с высокой вероятностью оказываются соседние точки сетки. Такая проверка может завышать качество, поскольку фактически оценивает локальную интерполяцию между близкими геометриями, а не устойчивость модели к исключению целого значения физического параметра. Для структурированных пространственных и параметрических данных необходимость более строгих схем валидации также обсуждается в литературе по structured cross-validation~\cite{roberts2017crossvalidation}.

Поэтому сравнение моделей проводится отдельно внутри каждого фиксированного класса $n$ с использованием двух протоколов. Первый протокол, LOAO, означает leave-one-$a$-out. Для каждого значения $a$ все геометрии с этим размером исключаются из обучения и используются как тестовая fold-группа:
\[
  \mathcal{D}_{\mathrm{test}}(a^\ast)
  =
  \{(a,r_{\mathrm{AR}}): a=a^\ast\}.
\]
Остальные значения $a$ образуют обучающую выборку. Такой протокол проверяет обобщение на невиданные значения размера внутри фиксированного класса формы.

Второй протокол, LOARO, означает leave-one-aspect-ratio-out. Для каждого значения $r_{\mathrm{AR}}$ из обучения исключаются все геометрии с данным отношением сторон:
\[
  \mathcal{D}_{\mathrm{test}}(r_{\mathrm{AR}}^\ast)
  =
  \{(a,r_{\mathrm{AR}}): r_{\mathrm{AR}}=r_{\mathrm{AR}}^\ast\}.
\]
Такой протокол проверяет обобщение на невиданные aspect ratios при сохранении фиксированного набора размеров.

Основной метрикой сравнения является MAE:
\begin{equation}
  \mathrm{MAE}
  =
  \frac{1}{M}
  \sum_{i=1}^{M}
  \left|\hat{y}_i-y_i\right| .
\end{equation}
Она вычисляется в физических единицах $t$. Это важно, поскольку суррогатная модель может использовать внутреннее масштабирование признаков или цели, но итоговая ошибка должна характеризовать реальную ошибку аппроксимации спектральной величины.

\section{Tiny MLP ablation}

Нелинейная модель в этой работе используется только как ablation/control experiment. Её задача --- проверить, даёт ли малая нелинейность устойчивое преимущество по сравнению с physics-informed Ridge при тех же структурированных протоколах валидации. MLP не рассматривается как источник физических законов и не заменяет прямой расчёт спектра в Kwant.

Используется намеренно малая архитектура. Для каждого фиксированного $n$ доступно только 35 точек параметрической сетки: пять значений $a$ и семь значений $r_{\mathrm{AR}}$. В таких условиях большая нейронная сеть легко может начать подгонять особенности конкретных fold-разбиений, особенно при наличии lattice-эффектов и малых зазоров между уровнями. Поэтому hidden layer состоит только из четырёх нейронов:
\[
  \texttt{hidden\_layer\_sizes=(4,)}.
\]
Вопрос о минимально необходимом объёме данных отдельно не исследовался, поэтому из результатов не делается утверждение о data efficiency.

Конфигурация tiny MLP фиксируется заранее:
\begin{itemize}
  \item внешний estimator: \texttt{TransformedTargetRegressor};
  \item масштабирование цели: \texttt{StandardScaler};
  \item regressor: \texttt{Pipeline};
  \item масштабирование признаков: \texttt{("scale", StandardScaler())};
  \item внутренняя модель: \texttt{MLPRegressor};
  \item \texttt{hidden\_layer\_sizes=(4,)};
  \item \texttt{activation="tanh"};
  \item \texttt{solver="lbfgs"};
  \item \texttt{alpha=1e-3};
  \item \texttt{early\_stopping=False};
  \item \texttt{random\_state=42} для фиксированной конфигурации;
  \item \texttt{max\_iter=5000}.
\end{itemize}

Для MLP рассматриваются два варианта признаков. Вариант MLP+raw использует исходные геометрические параметры $a$ и $r_{\mathrm{AR}}$. Вариант MLP+physics использует тот же physics-informed набор признаков, что и Ridge. Это разделение позволяет проверить две разные вещи: полезна ли малая нелинейность сама по себе и сохраняется ли преимущество физических признаков при нелинейной модели.

Поскольку обучение MLP может зависеть от начального состояния оптимизации, для оценки устойчивости используются несколько значений \texttt{random\_state}. Кроме того, для каждого fold фиксируется число итераций оптимизатора. Если модель достигает \texttt{max\_iter}, такой случай считается convergence failure и не может служить основанием для признания MLP устойчиво лучшей моделью.

\section{Предрегистрация критериев сравнения моделей}

Критерии сравнения Ridge и tiny MLP были зафиксированы до анализа результатов. Это необходимо, чтобы избежать post-hoc выбора модели по наиболее выгодным срезам таблицы ошибок. Сравнение проводится по ячейкам, определяемым тройкой:
\[
  (n,\; \mathrm{target},\; \mathrm{protocol}),
\]
где $n \in \{1.2,2.0,3.0,4.0\}$, target $\in \{E_0,dE_1\}$, а protocol $\in \{\mathrm{LOAO},\mathrm{LOARO}\}$. Всего получается
\[
  4 \times 2 \times 2 = 16
\]
основных comparison cells.

MLP+physics может считаться предпочтительной моделью только при устойчивом улучшении относительно Ridge+physics. Улучшение в каждой ячейке измеряется по MAE:
\begin{equation}
  \Delta_{\mathrm{MAE}}
  =
  \mathrm{MAE}_{\mathrm{Ridge}}
  -
  \overline{\mathrm{MAE}}_{\mathrm{MLP}},
\end{equation}
где среднее для MLP берётся по 10 значениям random seed. Относительное улучшение должно составлять не менее 15\%:
\begin{equation}
  \frac{\Delta_{\mathrm{MAE}}}{\mathrm{MAE}_{\mathrm{Ridge}}}
  \geq 0.15 .
\end{equation}
Порог 15\% выбран как консервативный практический минимум: меньшие отличия на малом структурированном датасете трудно отделить от fold-variability, seed-variability и дискретизационных эффектов границы.

Дополнительно проверяется seed stability. Для 10 запусков MLP определяется стандартная ошибка среднего:
\begin{equation}
  \mathrm{SE}
  =
  \frac{\mathrm{std}(\mathrm{MAE}_{\mathrm{MLP}})}{\sqrt{10}} .
\end{equation}
Улучшение считается устойчивым только если
\begin{equation}
  \Delta_{\mathrm{MAE}} > 0
  \quad \mathrm{and} \quad
  \Delta_{\mathrm{MAE}} > 2\,\mathrm{SE}.
\end{equation}
Таким образом, модель, которая в среднем хуже Ridge, не может быть признана стабильной даже при большой вариативности по seed.

Наконец, используется правило сходимости: если в данной comparison cell хотя бы один MLP+physics fit достигает \texttt{max\_iter}, эта ячейка не засчитывается как success независимо от значения MAE. Это правило предотвращает ситуацию, когда численно нестабильный результат интерпретируется как физически или статистически устойчивое улучшение.

Полное предпочтение MLP возможно только при выполнении критерия не только в сумме по всем 16 ячейкам, но и отдельно для обоих протоколов. Успех только в LOAO или только в LOARO рассматривается как частичное наблюдение, но не как основание заменить Ridge в качестве основной модели. Если предзарегистрированный критерий не выполняется, Ridge остаётся предпочтительной моделью благодаря простоте, интерпретируемости, стабильности и физической мотивированности признаков. Итог этого сравнения анализируется в главе 5.

\section{Выводы по главе 4}

В данной главе сформулирована методология суррогатного моделирования низкоэнергетического спектра модельных суперэллиптических квантовых точек. Основные выводы состоят в следующем.

\begin{enumerate}
  \item Набор признаков $\{1/a^2,\;1/(a^2 r_{\mathrm{AR}}),\;r_{\mathrm{AR}}\}$ выбран из физической картины низкоэнергетического scaling, подтверждённой в главе 3.
  \item Ridge-регрессия является основной интерпретируемой baseline-моделью: она линейна в физически мотивированных признаках, имеет малую ёмкость и стабилизируется $L_2$-регуляризацией.
  \item Случайное train/test-разбиение не используется как основной протокол, поскольку датасет является структурированной сеткой параметров. Вместо этого применяются LOAO и LOARO.
  \item Tiny MLP рассматривается только как контролируемая нелинейная ablation-модель, а не как источник физических законов или замена прямого Kwant-расчёта.
  \item Критерии сравнения моделей предзарегистрированы: учитываются порог улучшения MAE, устойчивость по seed, отсутствие convergence failures и выполнение требований отдельно для LOAO и LOARO.
  \item Численные результаты сравнения Ridge и tiny MLP рассматриваются в главе 5, где проверяется, выполняется ли предзарегистрированный критерий.
\end{enumerate}

\clearpage
